\section{completeness}

\begin{definition}[Cauchy sequences]
\mymark{***}
\index{sequence!Cauchy}
\index{Cauchy!Cauchy sequence}
Let $(X,d)$ be a metric space and $x_k$ a sequence of points in $X$.
We say that $x_k$ is a
\emph{Cauchy sequence}
\mymargin{Cauchy sequence}
if
\[
 \forall \eps>0\colon \exists n\in \NN\colon \forall j>n \colon \forall k > n \colon d(x_j,x_k) < \eps.
\]
\end{definition}

The property that defines Cauchy sequences
can also be written as:
\[
  \lim_{k \to +\infty} \sup_{j\ge k} d(x_j, x_k) = 0.
\]

\begin{theorem}[convergent sequences are Cauchy]
\mymark{**}%
\label{th:se_converge_cauchy}
Let $x_k\to x$ be a convergent sequence in a metric space $(X,d)$. Then $x_k$ is
Cauchy.
\end{theorem}
%
\begin{proof}
\mymark{**}
By definition, if $x_k \to x$ we have
\[
  \forall \eps>0\colon \exists n\in \NN \colon
  \forall k>n \colon d(x_k,x)< \eps.
\]
Applying the triangle inequality, for every $j,k>n$
we obtain the desired result:
\[
  d(x_j, x_k) \le d(x_k,x) + d(x,x_j) \le 2\eps.
\]
\end{proof}

\begin{definition}[completeness]
\mymark{***}
\mymargin{completeness}
\index{completeness}
A metric space $(X,d)$ is said to be \emph{complete}%
\mymargin{complete}\index{complete}
if every Cauchy sequence is convergent.
\end{definition}

The prototype of a complete metric space is $\RR$, as we will see in
theorem~\ref{th:R-completo}.
An example of a non-complete metric space is $\QQ$.
Indeed, if we take a
sequence $q_n \in\QQ$ with $q_n\to \sqrt 2$, the sequence $q_n$ is convergent
in $\RR$ and hence is Cauchy in $\RR$. But since $q_n\in \QQ$, it follows
that $q_n$ is Cauchy also in $\QQ$ (the Cauchy condition is the same).
However, in $\QQ$, such a sequence does not converge since $\sqrt 2\not \in
\QQ$.

\begin{definition}[Banach space]
A normed vector space is said to be a
\emph{Banach space}%
\mymargin{Banach space}%
\index{space!Banach}%
\index{Banach!Banach space}
if, as a metric space, it is complete.
If the norm is Euclidean, i.e., it derives from an inner product, the space is
called
a \emph{Hilbert space}.
\mymargin{Hilbert space}%
\index{space!Hilbert}
\end{definition}

\begin{lemma}
\label{lm:cauchy_limitata}
Every Cauchy sequence is bounded
(more precisely: if $x_n$ is a Cauchy sequence in a metric space $X$, then the
set $\ENCLOSE{x_n\colon n\in \NN}$
is a bounded set).
\end{lemma}
%
\begin{proof}
Let $x_k \in \RR$ be a Cauchy sequence.
Given $\eps>0$, we know there exists $N\in \NN$
such that for every $k,j>N$ we have
$d(x_k,x_j) < \eps$. 
In particular, for every $k>N$, we have
\[
  d(x_k, x_{N+1}) < \eps.
\]
Thus, chosen
\[
    R > \max\ENCLOSE{d(x_0,x_1), d(x_0,x_2), \dots, d(x_0,x_N), d(x_0, x_{N+1})
  + \eps}
\]
we observe that for every $k\in \NN$, $d(x_0, x_k)< R$ because 
if $k \le N$, we have chosen $R$ to be larger than $d(x_0,x_k)$, and if $k > N$,
then
\[
  d(x_0,x_k) \le d(x_0,x_{N+1}) + d(x_{N+1},x_k)
    \le d(x_0,x_{N+1}) + \eps < R.
\]

This means that for every $k\in \NN$, $x_k \in B_{R}(x_0)$ 
which is the definition of boundedness in a
metric space.
\end{proof}

\begin{lemma}
\label{lm:cauchy_estratta_convergente}
If a Cauchy sequence has a convergent subsequence, then the entire sequence is
convergent.
\end{lemma}
%
\begin{proof}
Let $x_k$ be the Cauchy sequence and let $x_{k_j}\to x$ be a convergent
subsequence.
Then for every $\eps>0$
there exists $m$ such that if $k,j>m$, then $d(x_k ,x_j) < \eps$.
Since $x_{k_j} \to x$, we can find $j$ such that $k_j > m$ and $d(x_{k_j},x) <
\eps$. But then
\[
  d(x_k,x) \le d(x_k, x_{k_j}) + d(x_{k_j},x)
   \le 2 \eps.
\]
And this is true for every $k > m$, from which it follows that the definition of
limit $x_k \to x$ is satisfied.
\end{proof}

\begin{theorem}[completeness of $\RR$]
\mymark{***}%
\mymargin{$\RR$ is complete}%
\index{completeness!of $\RR$}%
\label{th:R-completo}%
$\RR$ is complete.
\end{theorem}
%
\begin{proof}
\mymark{***}
We need to show that if $x_k$ is a Cauchy sequence in $\RR$, then $x_k$
converges.
By lemma~\ref{lm:cauchy_limitata}, we know that $x_k$ is bounded.
But then, by the Bolzano-Weierstrass theorem, we know that $x_k$ has
a convergent subsequence.
Thanks to lemma~\ref{lm:cauchy_estratta_convergente}, we can thus conclude
that the sequence $x_k$ itself is convergent.
\end{proof}

\begin{corollary}[completeness of $\RR^n$ and $\CC$]
The spaces $\RR^n$ and $\CC$ (with the usual Euclidean distance)
are complete.
\end{corollary}
%
\begin{proof}
It suffices to observe that the convergence (or the Cauchy condition) of a
sequence in $\RR^n$
occurs if and only if each component of the sequence is convergent (or Cauchy)
in $\RR$.
Thus, since $\RR$ is complete, $\RR^n$ is also complete. 
As a metric space, $\CC$ is isomorphic to $\RR^2$, hence it is also complete.
\end{proof}

\begin{theorem}[completeness of compact spaces]
Every compact metric space is complete.
\end{theorem}
%
\begin{proof}
Since the space is compact, every Cauchy sequence
admits a convergent subsequence.
But then, thanks to lemma~\ref{lm:cauchy_estratta_convergente},
the entire sequence converges, and thus the space is complete.
\end{proof}

\begin{theorem}[closed subsets in compact and complete spaces]%
\label{th:chiuso_completo}%
\label{th:compatto_completo}%
Let $(X,d)$ be a metric space and let $A\subset X$ be a closed subset in $X$. If
$X$ is compact, then $A$ is also compact; if $X$ is complete, then $A$ is also
complete.
\end{theorem}
\begin{proof}
If $X$ is compact, from every sequence in $A$, one can extract a convergent
subsequence to a point in $X$.
But since $A$ is closed, the point lies in $A$, and thus the subsequence is
convergent in $A$.

A Cauchy sequence in $A$ is also Cauchy in $X$. 
If $X$ is complete, such a sequence converges to a point in $X$. 
If $A$ is closed, that point is in $A$, and thus the sequence converges in $A$.
\end{proof}

\begin{theorem}[extendability of uniformly continuous functions]
\mymark{**}%
\label{th:estensione_uniformemente_continua}%
Let $f\colon A \subset X \to Y$ be a function defined on a subset $A$
of a metric space $X$ with values in a complete metric space $Y$
(for example, $X=\RR$ and $Y=\RR$). If $f$ is uniformly continuous, then
there exists a unique continuous function $\tilde f \colon \bar A \to Y$
such that for every $x\in A$, we have $\tilde f(x) = f(x)$.
Moreover, $\tilde f$ is also uniformly continuous.
\end{theorem}
%
\begin{proof}
The first observation is
that if a function $f$ is uniformly continuous, then $f$
maps Cauchy sequences to Cauchy sequences.
Indeed, if $f$ is uniformly continuous, we have
\[
 \forall \eps>0\colon \exists \delta>0 \colon d(x,y)<\delta \implies d(f(x),f(y))<\eps
\]
and if $a_n$ is Cauchy, we have
\[
 \exists N\in\NN \colon k,j>N \implies d(a_k,a_j)< \delta
\]
thus combining the two properties
we obtain the Cauchy condition for $f(a_k)$:
\[
\forall \eps>0\colon \exists N\in \NN\colon k,j>N \implies d(f(a_k),f(a_j))< \eps.
\]

Given a point $x \in \bar A$, there certainly exists $a_n \in A$ with $a_n\to
x$.
We want to show that $f(a_n)$ has a limit in $Y$ and that this limit does not
depend on the chosen sequence $a_n$. If $a_n \to a$, then $a_n$ is Cauchy in $X$
(theorem~\ref{th:se_converge_cauchy}).
From what was said before, we can assert that $f(a_n)$ is also Cauchy
and thus converges in $Y$, since $Y$ is complete by hypothesis.
Moreover, the $\lim f(a_n)$ does not depend on the chosen sequence because
if $b_n\to x$ were another sequence, we could construct a sequence
$c_n$ that alternates the points of $a_n$ and $b_n$ (setting, for example,
$c_{2n} = a_n$ and
$c_{2n+1} = b_n$). Also $c_n \to x$ and is Cauchy, thus $f(c_n)$ converges.
But $f(a_n)$ and $f(b_n)$ are subsequences of $f(c_n)$ and thus converge
to the same limit.
We have thus shown that for every $x\in \bar A$, there exists $\ell \in Y$
such that if $a_n\to x$, $a_n\in A$, then $f(a_n) \to \ell$. We can
therefore define $\tilde f(x)=\ell$. If there exists a continuous extension of
$f$, it can only be defined in this way, since continuous functions
map convergent sequences to convergent sequences.

On the other hand, we can show that $\tilde f$ is a uniformly continuous
function.
Indeed, for every $\eps>0$, we can choose $\delta>0$ given by the uniform
continuity of $f$: given $x,y\in \bar A$ with $d(x,y)<\delta/3$, there will
exist
$x_n,y_n \in A$ with $x_n\to x$, $y_n\to y$. We can then find $n$
sufficiently large such that
$d(x_n,x)<\delta/3$, $d(y_n,y)<\delta/3$,
$d(f(x_n),\tilde f(x)) < \eps$, $d(f(y_n),\tilde f(y)) < \eps$.
We will then have
$d(x_n,y_n)< d(x_n,x)+d(x,y)+d(y,y_n) < \delta$, thus $d(f(x_n),f(y_n))<\eps$
and in conclusion
\[
d(f(x),f(y))
\le d(f(x),f(x_n)) + d(f(x_n),f(y_n)) + d(f(y_n),f(y))
\le 3\eps.
\]
We have thus demonstrated the uniform continuity of $\tilde f$ with $3\eps$
instead of
$\eps$ and $\delta/3$
instead of $\delta$.
\end{proof}

\begin{definition}[Lipschitz]
\mymark{***}
Let $f\colon X \to Y$ be a function defined between two metric spaces.
Given $L\ge 0$,
we say that $f$ is $L$-Lipschitz if
\index{function!Lipschitz}
for every $x,y \in X$, we have
\[
  d(f(x),f(y)) \le L \cdot d(x,y).
\]
We say that $f$ is Lipschitz if there exists $L\ge 0$ such that $f$ is
$L$-Lipschitz.
\end{definition}

\begin{theorem}
\label{th:lipschitz_uniformemente_continua}%
\mymark{*}%
If $f\colon X \to Y$ is Lipschitz, then
$f$ is sequentially continuous, i.e.,
\[
  x_k \to x \implies f(x_k)\to f(x).
\]
\end{theorem}
%
\begin{proof}
If $x_k\to x$, it means that $d(x_k,x) \to 0$, therefore
\[
  d(f(x_k), f(x)) \le L \cdot d(x_k,x) \to 0.
\]
\end{proof}

We observe that the distance $d(x,y)$ of a metric space $X$ is always a
$1$-Lipschitz function with respect to each of the two variables $x$ and $y$.
Indeed, by the reverse triangle inequality, we have
\[
  \abs{d(x_1,y) - d(x_2,y)} \le d(x_1, x_2).
\]
Consequently, the norm of a normed space is also $1$-Lipschitz. 
In particular, the distance and the norm are continuous functions.

\begin{theorem}[Banach-Caccioppoli contraction principle or fixed point theorem]
\mymark{***}%
\mymargin{contraction theorem}%
\index{theorem!Banach-Caccioppoli}%
\index{theorem!contraction}%
\index{fixed point}%
\index{contraction}%
\index{Caccioppoli!contraction theorem}%
\index{Banach!contraction theorem}%
\label{th:banach-caccioppoli}%
Let $X$ be a non-empty complete metric space and let $f\colon X \to X$ be an
$L$-Lipschitz function with $L<1$ (we will say that $f$ is a
\emph{contraction}).
Then there exists a unique point
$x\in X$ such that $f(x) = x$.
\end{theorem}
%
\begin{proof}
\mymark{***}
Consider any point $p \in X$ and define
the sequence $x_k\in X$ through the recursive definition
\[
\begin{cases}
  x_0 = p \\
  x_{k+1} = f(x_k).
\end{cases}
\]
Since $f$ is $L$-Lipschitz, we will have
\begin{align*}
  d(x_2, x_1) &= d(f(x_1),f(x_0)) \le L \cdot d(x_1,x_0) \\
  d(x_3, x_2) &= d(f(x_2),f(x_1)) \le L \cdot d(x_2,x_1)
  \le L^2 \cdot d(x_1, x_0) \\
  d(x_4, x_3) &= d(f(x_3),f(x_2)) \le L \cdot d(x_3,x_2)
  \le L^3 \cdot d(x_1, x_0) \\
  &\vdots
\end{align*}
we can then prove inductively that
for every $m\in \NN$, we have
\[
  d(x_{m+1}, x_m) \le L^m \cdot d(x_1, x_0).
\]
But then for every $k\in \NN$ and for every $j>k$,
using the triangle inequality and summing the geometric series,
we have
\[
  d(x_k,x_j) \le \sum_{m=k}^{j-1} d(x_m, x_{m+1})
   \le \sum_{m=k}^{j-1} L^m \cdot d(x_1, x_0)
   = \frac{L^k-L^j}{1-L} d(x_1,x_0).
\]
Since $L<1$, if $k\to +\infty$ and $j>k$, this quantity tends to zero, and thus
it follows that $x_k$ is a Cauchy sequence. Being $X$ complete by hypothesis, we
know that the sequence converges $x_k \to x$ to a point $x\in X$.
By the continuity of $f$, taking the limit in the equation
$x_{k+1} = f(x_k)$, we obtain $x = f(x)$.
We have thus found a fixed point.
If $y\in X$ were another fixed point, we would have:
\[
  d(x,y) = d(f(x),f(y)) \le L \cdot d(x,y)
\]
which is absurd if $L<1$ and $x\neq y$.
\end{proof}